\chapter{Tổng kết và hướng phát triển}
Qua đề tài, nhóm em đã thực hiện các bước từ kiểm tra ảnh và nhãn, tiền xử lý, trích xuất đặc trưng đến huấn luyện, so sánh mô hình và phân tích lỗi. Từ các cấu hình đã thử, E6 có điểm validation cao nhất, còn đường học và kiểm định chéo cho thấy cần tiếp tục kiểm tra tính ổn định khi dữ liệu huấn luyện thay đổi.

Kết quả hiện tại vẫn có những giới hạn: số ảnh train của mỗi lớp ít, validation đã được dùng để chọn trong 56 cấu hình, một số ảnh trùng giữa các split và test V1 đã từng được quan sát. Ngoài ra, model E6 gốc hiện thiếu trong workspace và cần được khôi phục từ bản tin cậy khớp manifest nếu nhóm muốn chạy lại phần suy luận. Do đó, nhóm chưa thể khẳng định chất lượng triển khai thực tế hoặc đưa ra điểm test cuối cho E6.

Trong thời gian tới, nhóm em đề xuất kiểm tra kỹ các cặp ảnh trùng, bổ sung dữ liệu độc lập và thử các biến đổi ảnh có cơ sở từ phân tích lỗi. Trước khi đánh giá E6 trên test, nhóm cần được cho phép, đồng thời khóa model, seed, chính sách train và các chỉ số sẽ báo cáo. Kết quả cần được trình bày đầy đủ dù tốt hay xấu và không được dùng để tiếp tục chọn mô hình. Với những giới hạn hiện tại, hệ thống phù hợp làm tài liệu học tập và nghiên cứu ban đầu hơn là một công cụ nhận diện hoa đáng tin cậy trong mọi điều kiện thực tế.
\section{Kinh nghiệm rút ra}
Qua việc làm đề tài, nhóm em thấy rằng kết quả của một mô hình không chỉ phụ thuộc vào tên thuật toán mà còn phụ thuộc vào cách xử lý ảnh, đặc trưng được chọn, số lượng dữ liệu và cách chia tập. Một mô hình đơn giản như K-NN vẫn cần được kiểm tra kỹ: nếu chỉ nhìn điểm trên tập train, nhóm có thể đánh giá quá cao khả năng của mô hình đối với ảnh chưa gặp.

Nhóm cũng nhận thấy việc lưu cấu hình và dự đoán quan trọng không kém việc tạo biểu đồ. Khi có ID ảnh, nhãn và dự đoán, nhóm có thể tính lại một chỉ số hoặc kiểm tra các ảnh bị nhầm. Khi thiếu artifact cục bộ, người nhận dự án cần có lệnh chạy rõ ràng để tạo lại thay vì dựa vào một con số không biết nguồn gốc. Đây là lý do README và các notebook được tổ chức theo thứ tự chạy, bắt đầu từ dữ liệu chính thức.
\section{Các bước có thể thực hiện tiếp}
Trước mắt, các thành viên có thể kiểm tra lại dữ liệu và chạy notebook theo hướng dẫn, sau đó đối chiếu điểm validation với bảng kết quả đã ghi. Nếu xuất hiện khác biệt, cần kiểm tra phiên bản thư viện, seed, split và việc khớp ảnh với nhãn trước khi kết luận mô hình thay đổi. Các output sinh trong lần chạy mới cần được lưu riêng nếu nhóm muốn so sánh giữa các máy.

Về nghiên cứu tiếp theo, nhóm có thể xem các ảnh lỗi để đề xuất một thay đổi cụ thể về tiền xử lý hoặc đặc trưng, rồi đánh giá thay đổi đó bằng quy trình đã định trước trên train và validation. Nếu mục tiêu là báo cáo khả năng tổng quát hóa, cần ưu tiên dữ liệu độc lập, kiểm tra ảnh trùng và hạn chế số lần dùng tập đánh giá. Nhóm chưa thực hiện các bước đó trong báo cáo này và không đưa ra số liệu dự đoán cho chúng.
